\begin{trapbox}
	\begin{description}[style=nextline, leftmargin=2.8em, labelsep=0.5em, itemsep=0.35em]
		\item[\textbf{\textcolor{CTrap}{易错点一（期望频数计算）}}] 计算期望频数时直接用行合计或列合计代替，或忘记除以总频数。
		\item[\textbf{\textcolor{CTrap}{正确理解（公式记忆）}}] 期望频数必须为 $E_{ij}=R_iC_j/n$，即行合计乘列合计再除以总频数。
		\item[\textbf{\textcolor{CTrap}{错因分析（混淆概念）}}] 误以为期望频数就是行合计或列合计，忽略了联合概率的估计。
	\end{description}

	\medskip
	\begin{description}[style=nextline, leftmargin=2.8em, labelsep=0.5em, itemsep=0.35em]
		\item[\textbf{\textcolor{CTrap}{易错点二（自由度确定）}}] 计算自由度时误写为 $rc-1$ 或 $r+c-2$，导致查临界值错误。
		\item[\textbf{\textcolor{CTrap}{正确理解（自由含义）}}] 独立变量单元格数为 $(r-1)(c-1)$，扣除行列合计约束后的自由度数。
		\item[\textbf{\textcolor{CTrap}{错因分析（线性约束）}}] 未理解行列合计固定后实际可自由变化的单元格数，简单采用总单元格数减一。
	\end{description}

	\medskip
	\begin{description}[style=nextline, leftmargin=2.8em, labelsep=0.5em, itemsep=0.35em]
		\item[\textbf{\textcolor{CTrap}{易错点三（条件忽视）}}] 未检验期望频数条件或样本独立性，直接计算卡方并下结论。
		\item[\textbf{\textcolor{CTrap}{正确理解（必须验证）}}] 使用条件：各单元格期望频数 $\geq5$（或至少 $80\%$ 满足），且样本相互独立。
		\item[\textbf{\textcolor{CTrap}{错因分析（条件缺失）}}] 忽视前提导致卡方近似失效，结论可能不可靠。
	\end{description}
\end{trapbox}
